\scanpage{111}
donc aussi sur $\widetilde P_G^B$, image inverse du \inserted{$\mathbb U^B$-tors.} $P_B$ sur $M_G$.
On a donc un \deleted{\unsure{bon}} morphisme de topos
\diagram{\begin{tikzcd}[row sep=large,column sep=large]
\unsure{\mathfrak m|m}_{\vec C}\times\widetilde M_G^C \arrow[rr] \arrow[dr,dashed] && M_{G'} \\
& \widetilde M_G^C \arrow[ur,dashed] &
\end{tikzcd}}
s'insérant dans le diagramme cartésien
\diagram{\begin{tikzcd}[column sep=large,row sep=large]
\unsure{\mathfrak m|m}_{\vec C}\times\widetilde P_G^B \arrow[r] \arrow[d] & SP_{G'} \arrow[d] \\
\unsure{\mathfrak m|m}_{\vec C}\times\widetilde M_G^C \arrow[r] & M_{G'}
\end{tikzcd}}
Le \deleted{\unsure{bon}} morphisme $\unsure{\mu|m}_{\vec C}\times\widetilde M_G^C\longrightarrow M_{G'}$ joue le rôle de l'inclusion de $\unsure{\mathcal V|V}_{G,G'}^*$ \deleted{dans} (\unsure{trace} sur $M_{G'}$ du voisinage tubulaire de $M_G$ dans $\widehat M_{G'}$) dans $M_{G'}$.

Les choses \unsure{deviennent} plus claires peut-être (mais moins intrinsèques) si on dispose d'une section de $M_G^C$ sur $\unsure{M_{G_2}|M_G}$ (y en a-t-il toujours ??) auquel cas on trouve un diagramme cartésien
\margin{je pense que oui\ldots}
\diagram{\begin{tikzcd}[column sep=large,row sep=large]
SP_G\times m_C \arrow[r] \arrow[d] & SP_{G'} \arrow[d] \\
P_C\times m_C \arrow[r] & M_{G'}
\end{tikzcd}}
\mspage{110}
